\chapter{Implementation and Experimental Setup}
\label{chap:implementation}

\section{Development Environment}
\label{sec:dev_environment}
The CycleSafe computational environment is designed for containerized reproducibility across Unix and Windows hosts. The software stack is orchestrated via Docker Compose across five isolated microservices:

\begin{itemize}
    \item \textbf{Host Environment:} Windows 11 Enterprise / Ubuntu 22.04 LTS, Docker Engine v27.0 with Docker Compose v2.28.
    \item \textbf{Backend Core:} Go 1.23 standard library HTTP multiplexer, \texttt{pgx/v5} PostgreSQL driver, \texttt{go-redis/v9}.
    \item \textbf{Simulation Service:} Eclipse SUMO 1.27.1, Python 3.12, \texttt{traci} socket client, \texttt{xgboost} v2.0, \texttt{numpy}, \texttt{scipy}.
    \item \textbf{Database and In-Memory Cache:} PostgreSQL 16 on Debian Bookworm with relational tables for rides, points, incidents, and metadata; Redis 7 for asynchronous job queuing and lock coordination.
    \item \textbf{Frontend Interface:} Nginx alpine reverse proxy, Vanilla JavaScript (ES2022), Leaflet v1.9.4 with Canvas marker rendering, Chart.js v4.4.1.
\end{itemize}

\section{Hardware Telemetry Implementation}

\subsection{Component Procurement and Sourcing Specification}
\label{sec:hardware_bom}
To ensure practical feasibility within academic budget constraints, all components were sourced from domestic Indian electronics suppliers with verified shipping to \textbf{IIT (ISM) Dhanbad (PIN: 826004)}. Table \ref{tab:bom_table} details the complete Bill of Materials (BOM).

\begin{table}[h!]
	\centering
	\caption{Final verified Bill of Materials (BOM) for CycleSafe telemetry node.}
	\label{tab:bom_table}
	\setlength{\tabcolsep}{8pt}
	\renewcommand{\arraystretch}{1.25}
	\begin{adjustbox}{width=\textwidth}
	\begin{tabular}{|l|l|l|r|}
		\hline
		\textbf{Component Description} & \textbf{Specification / Model} & \textbf{Supplier (Delivers to Dhanbad 826004)} & \textbf{Cost (INR)} \\
		\hline
		\hline
		ESP32 Microcontroller (1 Unit) & DEVKIT V1 30-Pin CP2102 & Amazon India & 490 \\
		Ultrasonic Transducers (2 Units) & JSN-SR04T V3.0 Waterproof & Quartz Components & 590 \\
		GPS Module with Antenna (1 Unit) & u-blox NEO-6MV2 Ceramic & TechDelivers & 246 \\
		OLED Display Module (1 Unit) & 0.96" SSD1306 I2C 128x64 & Amazon India & 289 \\
		6-Axis IMU Module (1 Unit) & InvenSense MPU-6050 & Amazon India & 179 \\
		MicroSD SPI Reader (1 Unit) & Mini MicroSD 3.3V SPI & Quartz Components & 24 \\
		MicroSD Memory Card (1 Unit) & Geonix 32GB Class 10 & Amazon India & 479 \\
		Rechargeable Battery (1 Unit) & 18650 Li-ion 2500mAh 3.7V & Quartz Components & 67 \\
		Battery Holder (1 Unit) & Single 18650 with Leads & Amazon India & 45 \\
		Battery Charger Module (1 Unit) & TP4056 Type-C with Protection & Amazon India & 39 \\
		Voltage Regulator Module (1 Unit) & AMS1117-3.3V Step-Down & Quartz Components & 18 \\
		Momentary Push Button (1 Unit) & 12mm Waterproof Panel Mount & Amazon India & 228 \\
		Custom 2-Layer PCB (5 Boards) & OpenBikeSensor v00.03 FR-4 & LionCircuits Bangalore & 450 \\
		3D Printed Enclosure (1 Unit) & PETG Weatherproof Housing & Think3D India & 350 \\
		Lab Connectors \& Wire Harness & JST-XH 4-Pin, Standoffs, Solder & Department Lab Stock & 0 \\
		\hline
		\multicolumn{3}{|l|}{\textbf{Total Estimated Procurement Cost (INR)}} & \textbf{3,494} \\
		\hline
	\end{tabular}
	\end{adjustbox}
\end{table}

\subsection{Complete Microcontroller Pinout Allocation}
Table \ref{tab:pinout_table} documents the complete pin mapping on the 30-pin ESP32 DEVKIT V1 board.

\begin{table}[h!]
	\centering
	\caption{ESP32 30-pin microcontroller GPIO allocation and electrical bus mapping.}
	\label{tab:pinout_table}
	\setlength{\tabcolsep}{10pt}
	\renewcommand{\arraystretch}{1.25}
	\begin{tabular}{|l|l|l|l|}
		\hline
		\textbf{ESP32 Pin} & \textbf{Peripheral Device} & \textbf{Signal Function} & \textbf{Operating Rail} \\
		\hline
		\hline
		GPIO 21 & SSD1306 OLED + MPU-6050 & Shared I2C Data (SDA) & 3.3\,V (4.7\,k$\Omega$ Pullup) \\
		GPIO 22 & SSD1306 OLED + MPU-6050 & Shared I2C Clock (SCL) & 3.3\,V (4.7\,k$\Omega$ Pullup) \\
		GPIO 16 & u-blox NEO-6M GPS & Hardware UART2 RXD & 3.3\,V Logic \\
		GPIO 17 & u-blox NEO-6M GPS & Hardware UART2 TXD & 3.3\,V Logic \\
		GPIO 5  & MicroSD Breakout Board & SPI Chip Select (CS) & 3.3\,V Regulated \\
		GPIO 18 & MicroSD Breakout Board & SPI Clock (SCK) & 3.3\,V Regulated \\
		GPIO 19 & MicroSD Breakout Board & SPI Master-In Slave-Out (MISO) & 3.3\,V Regulated \\
		GPIO 23 & MicroSD Breakout Board & SPI Master-Out Slave-In (MOSI) & 3.3\,V Regulated \\
		GPIO 25 & Left JSN-SR04T Sensor & Ultrasonic Trigger Pulse & 3.3\,V Logic (5\,V VCC) \\
		GPIO 26 & Left JSN-SR04T Sensor & Ultrasonic Echo Duration & 3.3\,V Logic (Divider) \\
		GPIO 27 & Right JSN-SR04T Sensor & Ultrasonic Trigger Pulse & 3.3\,V Logic (5\,V VCC) \\
		GPIO 14 & Right JSN-SR04T Sensor & Ultrasonic Echo Duration & 3.3\,V Logic (Divider) \\
		GPIO 4  & Handlebar Push Button & Interrupt Input (Event Trigger) & Internal Pullup \\
		3V3     & Sensors, IMU, SD Card & Main Clean Voltage Rail & 3.3\,V LDO (AMS1117) \\
		GND     & All Subsystems & Common Ground Plane & 0\,V \\
		\hline
	\end{tabular}
\end{table}

\subsection{Mechanical Enclosure and 3D Print Parameters}
\label{sec:enclosure_specs}
The physical casing follows the OpenBikeSensor CAD specifications, divided into the \textbf{Main Sensor Enclosure} (seatpost mounted) and the \textbf{Handlebar Display Enclosure}. Both units are fabricated via Fused Deposition Modeling (FDM) in \textbf{PETG} filament. PETG is selected because its glass transition temperature ($T_g \approx 85^\circ\text{C}$) prevents softening under direct solar thermal radiation, unlike PLA ($T_g \approx 55^\circ\text{C}$), while offering high inter-layer fracture toughness against road vibration.

Production slicer settings include:
\begin{itemize}
    \item Layer Height: 0.20\,mm, Wall Loops: 4 perimeters (1.6\,mm solid boundary for water seal),
    \item Infill Density: 30\% Gyroid infill pattern, Extrusion Temperature: $240^\circ\text{C}$, Bed Temperature: $75^\circ\text{C}$.
\end{itemize}

\section{Experimental Simulation Setup}
\label{sec:simulation_setup}

\subsection{Corridor Alignment and OpenStreetMap Geometry}
The simulated corridor models the arterial thoroughfare connecting Charlottenburg and Friedrichshain-Kreuzberg in Berlin, Germany:
\begin{itemize}
    \item \textbf{Point A (Origin):} Ernst-Reuter-Platz ($52.51264^\circ\,\text{N}, 13.32189^\circ\,\text{E}$),
    \item \textbf{Point B (Destination):} Oberbaumbr{\"u}cke ($52.50120^\circ\,\text{N}, 13.44500^\circ\,\text{E}$),
    \item \textbf{Total Corridor Length:} 9,525.5\,m (~9.53\,km), discretized into 635 high-resolution geographic waypoints extracted from OpenStreetMap road alignments.
\end{itemize}

\subsection{Empirically Grounded Traffic Population}
Traffic injection rates are calibrated against the official Berlin Senate Oberbaumbr{\"u}cke counting station (3 April 2025, 07:00--19:00: 13,531 motor vehicles and 1,397 bicycles). Normalizing the 12-hour count to a standard 60-minute simulation hour yields the demand schedule in Table \ref{tab:sim_demand}.

\begin{table}[h!]
	\centering
	\caption{Stochastic vehicle population parameterized for 60-minute SUMO simulation.}
	\label{tab:sim_demand}
	\setlength{\tabcolsep}{10pt}
	\renewcommand{\arraystretch}{1.25}
	\begin{tabular}{|l|r|r|r|l|}
		\hline
		\textbf{Vehicle Class} & \textbf{12h Census} & \textbf{Hourly Mean} & \textbf{Simulated Fleet} & \textbf{SUMO Model Parameters} \\
		\hline
		\hline
		\textbf{Commuter Bicycles} & 1,397 & 116.4 & \textbf{115} & $v_{\text{max}} = 5.0\,\text{m/s}, a = 1.2, d = 3.0$ \\
		\textbf{Passenger Cars} & 12,260 & 1,021.7 & \textbf{1,024} & $v_{\text{max}} = 13.5\,\text{m/s}, a = 2.6, d = 4.5$ \\
		\textbf{Motorcycles (Krad)} & 595 & 49.6 & \textbf{52} & $v_{\text{max}} = 13.0\,\text{m/s}, a = 3.0, d = 5.0$ \\
		\textbf{Heavy Trucks \& Buses} & 676 & 56.3 & \textbf{53} & $v_{\text{max}} = 11.0\,\text{m/s}, a = 1.3, d = 4.0$ \\
		\hline
		\textbf{Total Fleet} & \textbf{14,928} & \textbf{1,244.0} & \textbf{1,244} & Total Simulated Duration: 3,600\,s \\
		\hline
	\end{tabular}
\end{table}

Vehicle departure timestamps $t_{\text{depart}} \in [0, 3600]$ are generated via a Poisson process. The simulation step size is set to $\Delta t = 0.5\,\text{s}$ (7,201 total simulation steps). Lane assignments adhere to German right-hand driving rules (\textit{Rechtsfahrgebot}): 65\% of motor vehicles depart in the primary mixed right lane, while 35\% depart in the left passing lane.

\section{Evaluation Metrics}
The quantitative performance of the classification engine and conflict detection pipeline is evaluated using standard machine learning metrics:
\begin{align}
    \text{Precision} &= \frac{\text{TP}}{\text{TP} + \text{FP}}, \quad
    \text{Recall} = \frac{\text{TP}}{\text{TP} + \text{FN}}, \quad
    F_1\text{-Score} = 2 \cdot \frac{\text{Precision} \cdot \text{Recall}}{\text{Precision} + \text{Recall}}
\end{align}
In the kinematic domain, conflict severity is benchmarked using Minimum Passing Distance ($d_{\text{min}}$), Peak Longitudinal Deceleration ($a_{\text{min}}$), and Time-to-Collision ($\text{TTC}$).

\section{Reproducibility Protocols}
All simulation configurations (\texttt{nodes.nod.xml}, \texttt{edges.edg.xml}, \texttt{corridor\_route.json}), Python TraCI orchestrators (\texttt{runner.py}, \texttt{predictor.py}), and trained model weights (\texttt{xgboost\_safety\_model.json}) are tracked under Git version control. By setting an explicit integer seed (e.g., \texttt{--seed 12345}), any experiment reported in this work can be reproduced identically with bit-level determinism.
